\begin{center}
    \textbf{Abstract}
\end{center}
%\phantomsection % required if using hyperref
%\addcontentsline{toc}{section}{\abstractname}
\doublespacing

Research on the impacts of mother tongue instruction (MTI) on labor market outcomes is limited. Point estimates suggest higher salary employment with MTI in Ethiopia, while wage-employment effects are imprecise. The salary-employment association is not statistically significant at 5\% using small-sample regional clustering adjustments; with only four regions, causal conclusions remain tentative. The positive salary-employment point estimate is consistent with the prevailing ethno-linguistic federalism in the country and a suggestive evidence in this study that cognitive ability increased with MTI. Being from a region with a historically dominant language does not confer advantages in the labor market. Institutional and policy measures may mediate the employment returns to MTI, but stronger evidence is needed.\\

\vspace{0in}
\noindent\textbf{Keywords:} Mother Tongue Instruction; Cognitive Ability; Employment;\\
\vspace{-.2in}\\
\noindent\textbf{JEL Codes:} I25, I28, J21, J24, J48\\

\bigskip

